% Zumdahl Chapter 4 -- Types of Chemical Reactions and Solution Stoichiometry.
% ELABORATED notes: more exposition, multiple worked examples per idea, and
% explicit reasoning chains. 5 block lessons.
\documentclass{apchem-notes}

\apcunitword{Chapter}
\apcunit{4}
\apcunittitle{Types of Chemical Reactions and Solution Stoichiometry}
\apcdoctitle{Guided Notes}
\apcsubtitle{Zumdahl \S4.1--4.11 \textbullet\ PDF pp.~168--225 \textbullet\ 5 blocks}

\begin{document}
\apctitleblock

\begin{apcnote}[How this chapter fits the AP course]
This is the closest the textbook comes to covering an entire CED unit in one
place: nearly all of \textbf{Unit 4 (Chemical Reactions)} lives here, and
\S4.1--4.3 backfill the solution and molarity content of \textbf{Unit
3.7--3.8}. Three skills introduced here are used continuously for the rest
of the year --- writing \emph{net ionic equations}, assigning
\emph{oxidation states}, and doing \emph{solution stoichiometry}, which
becomes titration math in Unit 8.

One caution: \S4.10's full half-reaction balancing in acidic and basic
solution goes beyond what the CED requires. It is marked
\textbf{ENRICHMENT} in Block 5 --- valuable, but not tested.
\end{apcnote}

% =====================================================================
\blocklesson{Water, Solutions, and Electrolytes}{Zumdahl \S4.1--4.3}

\begin{objectives}
  \item Explain why water dissolves ionic and polar substances, at the
        particle level.
  \item Classify solutes as strong electrolytes, weak electrolytes, or
        nonelectrolytes and predict conductivity.
  \item Calculate molarity, prepare solutions, and perform dilutions.
\end{objectives}

\reading{Zumdahl \S4.1--4.3}{169--182}

\begin{warmup}[3]
  \item Geometry and polarity of \ce{H2O}: \answerblank[4.5cm]{bent; polar}
  \item Formula for ammonium sulfate: \answerblank[3cm]{\ce{(NH4)2SO4}}
  \item Moles in \SI{58.44}{\gram} \ce{NaCl}: \answerblank[2.6cm]{\SI{1.000}{\mole}}
\end{warmup}

\segment{INSTRUCTION A}{25}{Why water is the universal solvent}

\notesectionz{The structure behind the behavior}{4.1}
\practice{1}

Water's power as a solvent is a direct consequence of two structural facts
you already know from Unit 2. First, each O--H bond is strongly
\answerblank[2.4cm]{polar}, because oxygen's electronegativity far exceeds
hydrogen's. Second, the molecule is \answerblank[1.6cm]{bent}, so those two
bond dipoles do \emph{not} cancel --- water carries a substantial net
dipole, with a partial negative charge on \answerblank[1.4cm]{oxygen} and
partial positives on the hydrogens.

That permanent dipole is the whole story. When an ionic solid is placed in
water, the water molecules orient themselves around each ion: oxygen ends
point toward \answerblank[2.4cm]{cations}, hydrogen ends point toward
\answerblank[2.4cm]{anions}. The resulting \term{ion--dipole} attractions
release enough energy to compete with the lattice energy holding the crystal
together. If ion--dipole attraction wins, the solid
\term{dissolves} and the ions become \term{hydrated} --- each surrounded by
a shell of oriented water molecules.

\begin{center}
\begin{tikzpicture}[font=\sffamily\footnotesize, scale=0.95]
  % central cation
  \node[circle, fill=black, text=white, inner sep=3pt] (ion) at (0,0) {$+$};
  \foreach \a in {0,60,...,300}{
    \begin{scope}[rotate=\a]
      \node[circle, draw, thick, inner sep=2.2pt] (o) at (1.45,0) {};
      \node[font=\tiny] at (1.45,0) {O};
      \draw[thick] (o) -- ++(28:0.42);
      \draw[thick] (o) -- ++(-28:0.42);
    \end{scope}
  }
  \node[below=1.9cm of ion, font=\sffamily\footnotesize]
    {oxygen ends turn toward the cation --- a hydrated ion};
\end{tikzpicture}
\end{center}

\begin{apcnote}
Dissolving is a \emph{competition}, not a guarantee. The lattice must be
pulled apart (costs energy) and the ions must be hydrated (releases energy).
When hydration cannot pay for the lattice, the salt stays solid --- which is
exactly what the solubility rules in Block 2 summarize empirically.
\end{apcnote}

\segment{GUIDED PRACTICE}{15}{Predicting from structure}

\begin{enumerate}[leftmargin=*,itemsep=0.35em]
  \item Which end of a water molecule approaches a \ce{Cl-} ion, and why?
        \answerlines[2]{The hydrogen end --- H carries $\delta^+$, which is
        attracted to the negative chloride ion.}
  \item Would you expect \ce{CCl4} to dissolve \ce{NaCl}? Justify.
        \answerlines[2]{No --- \ce{CCl4} is nonpolar and offers no
        ion--dipole attraction to compensate for \ce{NaCl}'s large lattice
        energy.}
\end{enumerate}

\segment{INSTRUCTION B}{20}{Electrolytes: three behaviors}

\notesectionz{Strong, weak, and non}{4.2}
\practice{1}

Conductivity is a \emph{measurement of ion concentration}. A solution
conducts only if mobile ions are present, so a conductivity test sorts
solutes into three classes:

\renewcommand{\arraystretch}{1.45}
\noindent\begin{tabular}{@{}p{3.1cm}p{4.3cm}p{5.9cm}@{}}
\toprule
\textbf{Class} & \textbf{What happens in water} & \textbf{Examples}\\
\midrule
\term{Strong electrolyte} &
  dissociates or ionizes \answerblank[2.2cm]{completely} &
  soluble salts (\ce{NaCl}); strong acids (\ce{HCl}, \ce{HNO3},
  \ce{H2SO4}); strong bases (\ce{NaOH}, \ce{KOH})\\
\term{Weak electrolyte} &
  ionizes only \answerblank[2.2cm]{slightly} --- an equilibrium &
  weak acids (\ce{HC2H3O2}); weak bases (\ce{NH3})\\
\term{Nonelectrolyte} &
  dissolves as \answerblank[3.4cm]{intact molecules} &
  sugar, ethanol; and \emph{water itself}\\
\bottomrule
\end{tabular}

\begin{workedexample}[Worked example: two acids, two behaviors]
\textbf{Hydrochloric acid} is a strong acid. In water:
\[ \ce{HCl(aq) -> H+(aq) + Cl-(aq)} \]
The single arrow means essentially every \ce{HCl} molecule has come apart.
A \SI{0.10}{\Molar} solution contains \SI{0.10}{\Molar} \ce{H+}.

\textbf{Acetic acid} is a weak acid. In water:
\[ \ce{HC2H3O2(aq) <=> H+(aq) + C2H3O2-(aq)} \]
The double arrow means the reaction reaches equilibrium with only about 1\%
of the molecules ionized. A \SI{0.10}{\Molar} solution contains roughly
\SI{0.001}{\Molar} \ce{H+} --- one hundred times less, and it lights the
conductivity bulb only dimly.
\end{workedexample}

\begin{apctrap}
``Strong'' and ``concentrated'' are unrelated words. \emph{Strong} describes
the \emph{fraction} that ionizes (a property of the substance);
\emph{concentrated} describes how much solute is present (a property of the
solution). A dilute \ce{HCl} solution is still a strong acid; a concentrated
acetic acid solution is still weak.
\end{apctrap}

\segment{APPLICATION}{20}{Molarity, dilution, and preparation}

\notesectionz{Concentration}{4.3}
\practice{5}

\[
  M = \answerblank[4.8cm]{$\dfrac{\text{moles of solute}}{\text{liters of solution}}$}
\]

Note carefully: liters of \emph{solution}, not liters of solvent. This is
why you dissolve the solid first and \emph{then} add water up to the mark.

\begin{workedexample}[Worked example: preparing a standard solution]
Prepare \SI{500.0}{\milli\liter} of \SI{0.250}{\Molar} \ce{Na2CO3}
($M = \SI{105.99}{\gram\per\mole}$).
\begin{align*}
  n &= (0.5000\ \text{L})(0.250\ \text{mol/L}) = \SI{0.125}{\mole}\\
  m &= (0.125)(105.99) = \SI{13.2}{\gram}
\end{align*}
Weigh \SI{13.2}{\gram}, dissolve in roughly \SI{300}{\milli\liter} of water,
transfer to a \SI{500}{\milli\liter} volumetric flask, then add water
\emph{to} the calibration mark and invert to mix.
\end{workedexample}

\begin{workedexample}[Worked example: dilution from a stock solution]
How much \SI{16}{\Molar} \ce{HNO3} is needed to make
\SI{250.}{\milli\liter} of \SI{0.50}{\Molar}?
\[
  V_1 = \frac{M_2V_2}{M_1} = \frac{(0.50)(250.)}{16} = \SI{7.8}{\milli\liter}
\]
Dilution changes volume but not \answerblank[3.2cm]{moles of solute}, which
is why $M_1V_1 = M_2V_2$ works. (Always add acid \emph{to} water, never the
reverse.)
\end{workedexample}

\notesub{Ion concentrations in electrolyte solutions}

Because strong electrolytes dissociate completely, the ion concentrations
follow directly from the formula:
\begin{enumerate}[leftmargin=*,itemsep=0.25em]
  \item \SI{0.10}{\Molar} \ce{Na2SO4}: $[\ce{Na+}] =$
        \answerblank[2cm]{\SI{0.20}{\Molar}}, $[\ce{SO4^2-}] =$
        \answerblank[2cm]{\SI{0.10}{\Molar}}
  \item \SI{0.25}{\Molar} \ce{Al(NO3)3}: $[\ce{NO3-}] =$
        \answerblank[2cm]{\SI{0.75}{\Molar}}
\end{enumerate}

\begin{exitticket}
A \SI{0.10}{\Molar} solution of substance X barely lights a conductivity
bulb; \SI{0.10}{\Molar} of substance Y lights it brightly. What can you
conclude about each, and what can you \emph{not} conclude?
\answerlines[3]{X is a weak electrolyte (only slightly ionized); Y is a
strong electrolyte (completely dissociated). You cannot conclude anything
about their concentrations --- both are 0.10 M --- nor about whether either
is an acid.}
\end{exitticket}

% =====================================================================
\blocklesson{Precipitation Reactions and Solubility}{Zumdahl \S4.4--4.5, 4.7}

\begin{objectives}
  \item Use solubility rules to predict whether a precipitate forms.
  \item Predict products by ion interchange and write balanced equations.
  \item Carry out stoichiometry for precipitation reactions.
\end{objectives}

\reading{Zumdahl \S4.4--4.5, 4.7}{182--192}

\begin{warmup}[4]
  \item $[\ce{Cl-}]$ in \SI{0.20}{\Molar} \ce{CaCl2}:
        \answerblank[2.6cm]{\SI{0.40}{\Molar}}
  \item Strong or weak electrolyte: \ce{HNO3}?
        \answerblank[2.6cm]{strong}
  \item Dilute \SI{10.0}{\milli\liter} of \SI{6.0}{\Molar} to
        \SI{60.0}{\milli\liter}: \answerblank[2.6cm]{\SI{1.0}{\Molar}}
  \item Nonelectrolyte example: \answerblank[3cm]{sugar / ethanol}
\end{warmup}

\segment{INSTRUCTION A}{25}{The solubility rules}

\notesectionz{Zumdahl's Table 4.1}{4.5}
\practice{4}

These rules are empirical --- they summarize what experiment shows, in
rough order of reliability. Learn them in order; the later rules are the
ones with the important exceptions.

\renewcommand{\arraystretch}{1.4}
\noindent\begin{tabular}{@{}p{0.5cm}p{6.2cm}p{6.4cm}@{}}
\toprule
 & \textbf{Soluble} & \textbf{Notable exceptions}\\
\midrule
1 & Most \answerblank[2.2cm]{nitrate} (\ce{NO3-}) salts & none worth
  memorizing\\
2 & Salts of alkali metals (\ce{Li+}, \ce{Na+}, \ce{K+}, \ce{Rb+},
  \ce{Cs+}) and \answerblank[2.2cm]{\ce{NH4+}} & none\\
3 & Most \ce{Cl-}, \ce{Br-}, \ce{I-} salts &
  \answerblank[3.9cm]{\ce{Ag+}, \ce{Pb^2+}, \ce{Hg2^2+}}\\
4 & Most \ce{SO4^2-} salts &
  \answerblank[6.6cm]{\ce{BaSO4}, \ce{PbSO4}, \ce{Hg2SO4}, \ce{CaSO4}}\\
5 & Most hydroxides are only \emph{slightly} soluble &
  soluble: \answerblank[2.6cm]{\ce{NaOH}, \ce{KOH}}; marginal:
  \ce{Ba(OH)2}, \ce{Sr(OH)2}, \ce{Ca(OH)2}\\
6 & Most \ce{S^2-}, \ce{CO3^2-}, \ce{CrO4^2-}, \ce{PO4^3-} salts are only
  slightly soluble & except those with the cations from
  \answerblank[1.4cm]{rule 2}\\
\bottomrule
\end{tabular}

\begin{apcnote}
Rules 1 and 2 are the two you can lean on hardest: \emph{anything} paired
with nitrate, an alkali metal, or ammonium dissolves. When scanning a
possible product, check those two first --- they resolve most questions
instantly.
\end{apcnote}

\segment{GUIDED PRACTICE}{15}{Soluble or not?}

\begin{enumerate}[leftmargin=*,itemsep=0.3em]
  \item \ce{KNO3}: \answerblank[4.5cm]{soluble (rules 1 and 2)}
  \item \ce{AgCl}: \answerblank[5.6cm]{insoluble (rule 3 exception)}
  \item \ce{BaSO4}: \answerblank[5.6cm]{insoluble (rule 4 exception)}
  \item \ce{(NH4)2S}: \answerblank[5.7cm]{soluble (rule 2 beats rule 6)}
  \item \ce{CaCO3}: \answerblank[4.5cm]{insoluble (rule 6)}
\end{enumerate}

\segment{INSTRUCTION B}{20}{Predicting products by ion interchange}

\notesectionz{What forms when solutions are mixed?}{4.5}
\practice{4}

Zumdahl is candid that predicting products is one of the hardest things a
beginning student is asked to do. The procedure that makes it tractable:

\begin{enumerate}[leftmargin=*,itemsep=0.2em]
  \item List the ions actually present after mixing (remember: strong
        electrolytes are already \answerblank[2.2cm]{dissociated}).
  \item Form the two possible new pairings by
        \answerblank[3cm]{interchanging} cation and anion. A solid must be
        electrically neutral, so pair cation with
        \answerblank[1.6cm]{anion} only.
  \item Check each possibility against the solubility rules.
  \item Anything insoluble is the \term{precipitate}; the rest stays
        dissolved.
\end{enumerate}

\begin{workedexample}[Worked example: Zumdahl's chromate reaction]
Mix \ce{K2CrO4(aq)} and \ce{Ba(NO3)2(aq)}. Ions present:
\ce{K+}, \ce{CrO4^2-}, \ce{Ba^2+}, \ce{NO3-}.

Interchange gives two candidates: \ce{KNO3} and \ce{BaCrO4}.
Rules 1 and 2 make \ce{KNO3} soluble. Rule 6 makes \ce{BaCrO4} insoluble
(barium is not an alkali metal). So a yellow solid of \ce{BaCrO4} forms:
\[ \ce{K2CrO4(aq) + Ba(NO3)2(aq) -> BaCrO4(s) + 2 KNO3(aq)} \]
The \ce{K+} and \ce{NO3-} ions stay dispersed in solution. Remove the solid
and evaporate the water, and they would reassemble into solid \ce{KNO3}.
\end{workedexample}

\begin{workedexample}[Worked example: no reaction]
Mix \ce{NaCl(aq)} and \ce{KNO3(aq)}. Candidates: \ce{NaNO3} and \ce{KCl}.
Both are soluble by rules 1--3. \textbf{No precipitate forms} --- the
solution simply contains four kinds of ion. Write ``no reaction.'' Being
willing to answer ``nothing happens'' is part of the skill.
\end{workedexample}

\segment{APPLICATION}{20}{Precipitation stoichiometry}

The mass of precipitate is an ordinary stoichiometry problem --- with one
extra front step: convert volume and molarity into
\answerblank[1.6cm]{moles}.

\begin{workedexample}[Worked example]
What mass of \ce{AgCl} forms when \SI{50.0}{\milli\liter} of
\SI{0.200}{\Molar} \ce{AgNO3} is mixed with excess \ce{NaCl}?
\begin{align*}
  n(\ce{Ag+}) &= (0.0500)(0.200) = \SI{0.0100}{\mole}\\
  &\ce{Ag+ + Cl- -> AgCl(s)} \quad (1{:}1)\\
  m(\ce{AgCl}) &= (0.0100)(143.32) = \SI{1.43}{\gram}
\end{align*}
\end{workedexample}

\begin{enumerate}[leftmargin=*,itemsep=0.4em]
  \item \SI{25.0}{\milli\liter} of \SI{0.100}{\Molar} \ce{BaCl2} is mixed
        with excess \ce{Na2SO4}. Find the mass of \ce{BaSO4}
        ($M = 233.39$). \workspace[2.4cm]
        \keyonly{\begin{apcanswer}$n = (0.0250)(0.100) = \SI{2.50e-3}{\mole}$;
        1:1 $\to (2.50\times10^{-3})(233.39) =
        \SI{0.583}{\gram}$\end{apcanswer}}
  \item Explain why the answer does not depend on how much \ce{Na2SO4} was
        added, provided it is in excess. \answerlines[2]{\ce{BaCl2} is the
        limiting reactant; once all the \ce{Ba^2+} has precipitated, extra
        sulfate has nothing left to react with.}
\end{enumerate}

\begin{exitticket}
Predict whether mixing \ce{Pb(NO3)2(aq)} and \ce{KI(aq)} produces a
precipitate. Name it and cite the rule.
\answerlines[2]{Yes --- \ce{PbI2} is insoluble (rule 3 exception for
\ce{Pb^2+}); \ce{KNO3} stays dissolved by rules 1 and 2.}
\end{exitticket}

% =====================================================================
\blocklesson{Ionic Equations}{Zumdahl \S4.6}

\begin{objectives}
  \item Write formula, complete ionic, and net ionic equations.
  \item Identify spectator ions and explain why they are omitted.
\end{objectives}

\reading{Zumdahl \S4.6}{188--190}

\begin{warmup}[3]
  \item Soluble or not: \ce{PbSO4}? \answerblank[3cm]{insoluble}
  \item Products of \ce{AgNO3(aq) + KCl(aq)}:
        \answerblank[4.5cm]{\ce{AgCl(s)} + \ce{KNO3(aq)}}
  \item Moles of \ce{Ag+} in \SI{20.0}{\milli\liter} of
        \SI{0.150}{\Molar} \ce{AgNO3}:
        \answerblank[3cm]{\SI{3.00e-3}{\mole}}
\end{warmup}

\segment{INSTRUCTION A}{25}{Three equations, three levels of honesty}

\notesectionz{From formula to net ionic}{4.6}
\practice{1}

The formula equation is convenient but, as Zumdahl puts it, ``does not give
a correct picture of what actually occurs in solution.'' Three levels:

\renewcommand{\arraystretch}{1.4}
\noindent\begin{tabular}{@{}p{3.4cm}p{10.2cm}@{}}
\toprule
\textbf{Formula equation} & every substance written as a neutral compound\\
\textbf{Complete ionic} & every \answerblank[3.4cm]{strong electrolyte}
  written as separated ions\\
\textbf{Net ionic} & only the species that
  \answerblank[3cm]{actually react}\\
\bottomrule
\end{tabular}

\begin{workedexample}[Worked example: the same reaction, three ways]
\textbf{Formula:}
\[ \ce{K2CrO4(aq) + Ba(NO3)2(aq) -> BaCrO4(s) + 2 KNO3(aq)} \]
\textbf{Complete ionic} --- dissolve every strong electrolyte:
\[ \ce{2 K+(aq) + CrO4^2-(aq) + Ba^2+(aq) + 2 NO3^-(aq) -> BaCrO4(s) + 2 K+(aq) + 2 NO3^-(aq)} \]
\ce{K+} and \ce{NO3-} appear unchanged on both sides. They are
\term{spectator ions}. Cancel them:\\
\textbf{Net ionic:}
\[ \ce{Ba^2+(aq) + CrO4^2-(aq) -> BaCrO4(s)} \]
\end{workedexample}

\notesub{The rule that governs everything: what gets split?}

\begin{itemize}[leftmargin=*,itemsep=0.25em]
  \item \textbf{Split} (write as ions): soluble salts, strong acids, strong
        bases --- i.e.\ all \answerblank[3.9cm]{strong electrolytes}.
  \item \textbf{Keep together}: solids \ce{(s)}, liquids \ce{(l)}, gases
        \ce{(g)}, \answerblank[3.6cm]{weak electrolytes}, and
        nonelectrolytes.
\end{itemize}

\begin{apctrap}
Water is a \emph{nonelectrolyte} --- never write it as \ce{H+ + OH-} in an
ionic equation. And a weak acid such as \ce{HC2H3O2} stays intact even
though it is aqueous, because it barely ionizes. Splitting a weak acid is
the most common net-ionic error.
\end{apctrap}

\segment{GUIDED PRACTICE}{15}{Write all three}

For \ce{Pb(NO3)2(aq) + 2 NaI(aq) -> PbI2(s) + 2 NaNO3(aq)}:
\begin{enumerate}[leftmargin=*,itemsep=0.35em]
  \item Complete ionic: \answerlines[2]{\ce{Pb^2+(aq) + 2 NO3^-(aq) +
        2 Na+(aq) + 2 I^-(aq) -> PbI2(s) + 2 Na+(aq) + 2 NO3^-(aq)}}
  \item Spectator ions: \answerblank[4cm]{\ce{Na+} and \ce{NO3-}}
  \item Net ionic: \answerblank[6.6cm]{\ce{Pb^2+(aq) + 2 I^-(aq) -> PbI2(s)}}
\end{enumerate}

\segment{INSTRUCTION B}{20}{Why net ionic equations matter}

\notesectionz{The same reaction, many disguises}{4.6}
\practice{6}

The net ionic equation strips away everything incidental, which reveals that
apparently different reactions are the \emph{same} reaction.

\begin{workedexample}[Worked example: one net ionic, three formula equations]
All three of these produce the identical net ionic equation
\ce{Ag+(aq) + Cl^-(aq) -> AgCl(s)}:
\begin{align*}
  \ce{AgNO3(aq) + NaCl(aq) &-> AgCl(s) + NaNO3(aq)}\\
  \ce{AgNO3(aq) + KCl(aq) &-> AgCl(s) + KNO3(aq)}\\
  \ce{AgC2H3O2(aq) + NH4Cl(aq) &-> AgCl(s) + NH4C2H3O2(aq)}
\end{align*}
The spectators differ; the chemistry does not. This is why a chemist writes
the net ionic equation --- it names the actual event.
\end{workedexample}

\segment{APPLICATION}{20}{Net ionic reasoning}

\begin{enumerate}[leftmargin=*,itemsep=0.4em]
  \item Write the net ionic equation for mixing \ce{Na2CO3(aq)} and
        \ce{CaCl2(aq)}. \answerlines[2]{\ce{Ca^2+(aq) + CO3^2-(aq) ->
        CaCO3(s)} --- \ce{Na+} and \ce{Cl-} are spectators.}
  \item Mixing \ce{NaNO3(aq)} and \ce{KCl(aq)} gives no net ionic equation
        at all. Explain what that means physically.
        \answerlines[2]{Every possible product is soluble, so no ions
        combine --- all four ion types remain dispersed. Nothing has
        reacted.}
  \item A student writes \ce{Ba^2+ + SO4^2- -> BaSO4(s)} for the reaction of
        \ce{Ba(OH)2(aq)} with \ce{H2SO4(aq)}. What did they miss?
        \answerlines[3]{Two reactions occur: the sulfate precipitates AND
        \ce{H+} reacts with \ce{OH-} to form water. The full net ionic is
        \ce{Ba^2+ + 2 OH^- + 2 H+ + SO4^2- -> BaSO4(s) + 2 H2O(l)}.}
\end{enumerate}

\begin{exitticket}
Why is \ce{HC2H3O2(aq)} written as a whole molecule in a net ionic equation
while \ce{HCl(aq)} is written as separated ions?
\answerlines[2]{Acetic acid is a weak electrolyte --- only about 1\% ionizes,
so the molecular form dominates. \ce{HCl} is a strong acid and ionizes
essentially completely.}
\end{exitticket}

% =====================================================================
\blocklesson{Acid--Base Reactions}{Zumdahl \S4.8}

\begin{objectives}
  \item Define acids and bases in Arrhenius and Bronsted--Lowry terms.
  \item Write net ionic equations for strong and weak acid neutralizations.
  \item Perform titration stoichiometry.
\end{objectives}

\reading{Zumdahl \S4.8}{192--199}

\begin{warmup}[4]
  \item Net ionic for \ce{AgNO3(aq) + NaCl(aq)}:
        \answerblank[5.5cm]{\ce{Ag+ + Cl^- -> AgCl(s)}}
  \item Do you split \ce{H2O} in an ionic equation?
        \answerblank[2cm]{no}
  \item Strong acid example: \answerblank[4.1cm]{\ce{HCl}, \ce{HNO3},
        \ce{H2SO4}}
  \item Moles in \SI{35.0}{\milli\liter} of \SI{0.200}{\Molar}:
        \answerblank[3cm]{\SI{7.00e-3}{\mole}}
\end{warmup}

\segment{INSTRUCTION A}{25}{Two definitions, one behavior}

\notesectionz{Arrhenius and Bronsted--Lowry}{4.8}
\practice{1}

\begin{itemize}[leftmargin=*,itemsep=0.25em]
  \item \term{Arrhenius}: an acid produces \answerblank[1.6cm]{\ce{H+}} in
        water; a base produces \answerblank[2cm]{\ce{OH-}}.
  \item \term{Bronsted--Lowry}: an acid is a proton
        \answerblank[1.8cm]{donor}; a base is a proton
        \answerblank[1.8cm]{acceptor}. This is broader --- it covers
        \ce{NH3}, which produces \ce{OH-} without containing it.
\end{itemize}

\notesub{The strong acids and bases worth memorizing}

Strong acids: \ce{HCl}, \ce{HBr}, \ce{HI}, \ce{HNO3}, \ce{H2SO4},
\ce{HClO4}. \\
Strong bases: hydroxides of \answerblank[3.6cm]{group 1} plus
\ce{Ca(OH)2}, \ce{Sr(OH)2}, \ce{Ba(OH)2}. \\
Everything else you meet this year is \answerblank[1.6cm]{weak}.

\segment{GUIDED PRACTICE}{15}{Classify}

\begin{enumerate}[leftmargin=*,itemsep=0.3em]
  \item \ce{HNO3}: \answerblank[3.4cm]{strong acid}
  \item \ce{NH3}: \answerblank[3.4cm]{weak base}
  \item \ce{HF}: \answerblank[3.4cm]{weak acid}
  \item \ce{KOH}: \answerblank[3.4cm]{strong base}
\end{enumerate}

\segment{INSTRUCTION B}{20}{Neutralization at the ion level}

\notesectionz{Two cases, two net ionic equations}{4.8}
\practice{6}

\begin{workedexample}[Case 1: strong acid + strong base]
Mix \ce{HCl(aq)} and \ce{NaOH(aq)}. The mixed solution contains
\ce{H+}, \ce{Cl-}, \ce{Na+}, \ce{OH-}. Would \ce{NaCl} precipitate? Rules 2
and 3 say no --- so \ce{Na+} and \ce{Cl-} are spectators. But \ce{H+} and
\ce{OH-} cannot coexist in quantity, because water is a nonelectrolyte:
\[ \ce{H+(aq) + OH^-(aq) -> H2O(l)} \]
\textbf{Every} strong-acid/strong-base neutralization has this same net
ionic equation.
\end{workedexample}

\begin{workedexample}[Case 2: weak acid + strong base]
Mix \ce{HC2H3O2(aq)} and \ce{KOH(aq)}. Acetic acid is a weak electrolyte,
so it stays intact:
\[ \ce{HC2H3O2(aq) + OH^-(aq) -> C2H3O2^-(aq) + H2O(l)} \]
The whole acid molecule appears, because the hydrogen is transferred
directly rather than being freely available beforehand. Only \ce{K+} is a
spectator.
\end{workedexample}

\begin{apctrap}
The two cases give \emph{different} net ionic equations, and telling them
apart requires knowing whether the acid is strong or weak. Memorize the six
strong acids --- everything else is weak, and stays whole.
\end{apctrap}

\segment{APPLICATION}{20}{Titration stoichiometry}

A \term{titration} adds a solution of known concentration (the
\term{titrant}) until the reaction is exactly complete --- the
\term{equivalence point}, signaled by an
\answerblank[2.6cm]{indicator} colour change.

\begin{workedexample}[Worked example]
It takes \SI{27.4}{\milli\liter} of \SI{0.150}{\Molar} \ce{NaOH} to
neutralize \SI{25.0}{\milli\liter} of \ce{HCl}. Find $[\ce{HCl}]$.
\begin{align*}
  n(\ce{OH-}) &= (0.0274)(0.150) = \SI{4.11e-3}{\mole}\\
  &\ce{H+ + OH^- -> H2O} \quad (1{:}1) \Rightarrow n(\ce{H+}) = \SI{4.11e-3}{\mole}\\
  [\ce{HCl}] &= \frac{4.11\times10^{-3}}{0.0250} = \SI{0.164}{\Molar}
\end{align*}
\end{workedexample}

\begin{enumerate}[leftmargin=*,itemsep=0.4em]
  \item \SI{20.0}{\milli\liter} of \ce{H2SO4} requires
        \SI{32.0}{\milli\liter} of \SI{0.250}{\Molar} \ce{NaOH}. Find
        $[\ce{H2SO4}]$. Careful with the ratio. \workspace[2.6cm]
        \keyonly{\begin{apcanswer}$n(\ce{OH-}) = (0.0320)(0.250) =
        \SI{8.00e-3}{\mole}$. \ce{H2SO4} supplies 2 \ce{H+} per formula
        unit, so $n(\ce{H2SO4}) = 4.00\times10^{-3}$ mol;
        $[\ce{H2SO4}] = 4.00\times10^{-3}/0.0200 =
        \SI{0.200}{\Molar}$\end{apcanswer}}
  \item Why does the diprotic acid change the arithmetic?
        \answerlines[2]{Each \ce{H2SO4} donates two protons, so it takes
        twice as much base per mole of acid --- the mole ratio is 1:2, not
        1:1.}
\end{enumerate}

\begin{exitticket}
Write the net ionic equation for \ce{HF(aq)} reacting with \ce{NaOH(aq)},
and explain why it differs from the \ce{HCl}/\ce{NaOH} case.
\answerlines[2]{\ce{HF(aq) + OH^-(aq) -> F^-(aq) + H2O(l)}. HF is a weak
acid, so it stays intact rather than appearing as \ce{H+}.}
\end{exitticket}

% =====================================================================
\blocklesson{Oxidation--Reduction Reactions}{Zumdahl \S4.9--4.11}

\begin{objectives}
  \item Assign oxidation states using the standard rules.
  \item Identify what is oxidized, what is reduced, and name the agents.
  \item Balance simple redox equations; perform redox titration
        stoichiometry.
\end{objectives}

\reading{Zumdahl \S4.9--4.11}{199--212}

\begin{warmup}[3]
  \item Net ionic for strong acid + strong base:
        \answerblank[5cm]{\ce{H+ + OH^- -> H2O}}
  \item Moles of \ce{H+} in \SI{25.0}{\milli\liter} of \SI{0.100}{\Molar}
        \ce{H2SO4}: \answerblank[3cm]{\SI{5.00e-3}{\mole}}
  \item Charge on the sulfate ion: \answerblank[2cm]{$2-$}
\end{warmup}

\segment{INSTRUCTION A}{25}{Oxidation states}

\notesectionz{Bookkeeping for electrons}{4.9}
\practice{5}

An \term{oxidation state} is the \emph{imaginary} charge an atom would carry
if every shared pair were awarded entirely to the more electronegative
partner. It is a bookkeeping device, not a real charge --- but it lets us
track where electrons go.

\renewcommand{\arraystretch}{1.35}
\noindent\begin{tabular}{@{}p{0.5cm}p{6.6cm}p{6.0cm}@{}}
\toprule
 & \textbf{Rule} & \textbf{Note}\\
\midrule
1 & Free element: \answerblank[1.4cm]{0} & \ce{Na}, \ce{O2}, \ce{P4}\\
2 & Monatomic ion: equal to its \answerblank[1.6cm]{charge} &
  \ce{Na+} is $+1$\\
3 & Oxygen: \answerblank[1.4cm]{$-2$} & except peroxides ($-1$)\\
4 & Hydrogen: \answerblank[1.4cm]{$+1$} & except metal hydrides ($-1$)\\
5 & Fluorine: \answerblank[1.4cm]{$-1$} & always\\
6 & Sum $=$ \answerblank[1.2cm]{0} for a neutral compound, or the
  \answerblank[2.4cm]{ion charge} for an ion & this is the equation you
  actually solve\\
\bottomrule
\end{tabular}

\begin{workedexample}[Worked example: three species, Zumdahl's own]
\textbf{\ce{CO2}}: assign O first at $-2$ each ($-4$ total). The molecule is
neutral, so C must be $+4$. \emph{Reality check:} $(+4) + 2(-2) = 0$.

\textbf{\ce{SF6}}: no rule for sulfur, so start with F at $-1$ each ($-6$
total). S must be $+6$. Check: $(+6) + 6(-1) = 0$.

\textbf{\ce{NO3-}}: three O at $-2$ gives $-6$. The ion's charge is $-1$, so
N must be $+5$. Check: $(+5) + 3(-2) = -1$. \checkmark
\end{workedexample}

\begin{apcnote}
Non-integer oxidation states can occur. In \ce{Fe3O4}, four oxygens give
$-8$, so three irons must total $+8$ --- each is $+8/3$. That looks strange,
but the rules deliberately treat all three iron atoms as equivalent when the
compound is really $2\,\ce{Fe^3+} + 1\,\ce{Fe^2+}$ per formula unit.
\end{apcnote}

\segment{GUIDED PRACTICE}{15}{Assign the underlined element}

\begin{enumerate}[leftmargin=*,itemsep=0.3em]
  \item S in \ce{H2SO4}: \answerblank[2cm]{$+6$}
  \item Cr in \ce{Cr2O7^2-}: \answerblank[2cm]{$+6$}
  \item Mn in \ce{MnO4-}: \answerblank[2cm]{$+7$}
  \item N in \ce{NH3}: \answerblank[2cm]{$-3$}
  \item C in \ce{CH4}: \answerblank[2cm]{$-4$}
\end{enumerate}

\segment{INSTRUCTION B}{20}{Identifying redox and naming agents}

\notesectionz{Who lost, who gained}{4.9}
\practice{6}

A reaction is \term{redox} if any element's oxidation state
\answerblank[2cm]{changes}.

\renewcommand{\arraystretch}{1.3}
\noindent\begin{tabular}{@{}p{3.2cm}p{3.6cm}p{6.4cm}@{}}
\toprule
 & \textbf{Oxidation} & \textbf{Reduction}\\
\midrule
Electrons & \answerblank[1.4cm]{lost} & \answerblank[1.4cm]{gained}\\
Oxidation state & \answerblank[1.8cm]{increases} &
  \answerblank[1.8cm]{decreases}\\
The species itself is the & \answerblank[3.2cm]{reducing agent} &
  \answerblank[3.2cm]{oxidizing agent}\\
\bottomrule
\end{tabular}

\begin{apctrap}
The agent labels are counterintuitive by design: the substance that
\emph{is oxidized} is the \emph{reducing} agent, because it hands its
electrons to something else and thereby reduces it. Read the label as
``the thing that causes reduction.'' \textbf{OIL RIG} --- Oxidation Is Loss,
Reduction Is Gain.
\end{apctrap}

\begin{workedexample}[Worked example: full analysis]
\[ \ce{2 Al(s) + 3 Cl2(g) -> 2 AlCl3(s)} \]
Al goes $0 \to +3$: it \textbf{lost} electrons, so it was \textbf{oxidized}
and is the \textbf{reducing agent}.
Cl goes $0 \to -1$: it \textbf{gained} electrons, so it was
\textbf{reduced} and is the \textbf{oxidizing agent}.
Both reactants are free elements at $0$ --- a reliable sign that a reaction
is redox.
\end{workedexample}

\segment{APPLICATION}{20}{Redox titration}

\begin{enumerate}[leftmargin=*,itemsep=0.4em]
  \item For \ce{Zn(s) + Cu^2+(aq) -> Zn^2+(aq) + Cu(s)}, identify what is
        oxidized, reduced, and both agents. \answerlines[3]{Zn: $0 \to +2$,
        oxidized, reducing agent. Cu: $+2 \to 0$, reduced, oxidizing agent.}
  \item In an iron analysis, \ce{Fe^2+} is titrated with \ce{MnO4-}:
        \[ \ce{5 Fe^2+ + MnO4^- + 8 H+ -> 5 Fe^3+ + Mn^2+ + 4 H2O} \]
        A sample requires \SI{18.0}{\milli\liter} of \SI{0.0200}{\Molar}
        \ce{KMnO4}. Find the moles of \ce{Fe^2+}. \workspace[2.4cm]
        \keyonly{\begin{apcanswer}$n(\ce{MnO4-}) = (0.0180)(0.0200) =
        \SI{3.60e-4}{\mole}$; ratio 5:1 $\to n(\ce{Fe^2+}) =
        \SI{1.80e-3}{\mole}$\end{apcanswer}}
  \item Verify that manganese is reduced in that equation by computing its
        oxidation state on both sides.
        \answerlines[2]{In \ce{MnO4-}: $x + 4(-2) = -1 \Rightarrow x = +7$.
        As \ce{Mn^2+}: $+2$. It decreased, so Mn was reduced.}
\end{enumerate}

\begin{apcnote}[ENRICHMENT --- beyond the CED]
Zumdahl \S4.10 develops full \emph{half-reaction balancing} in acidic and
basic solution, adding \ce{H2O}, \ce{H+}, and \ce{OH-} to balance oxygen and
hydrogen. The current AP framework does not require this --- it asks you to
\emph{identify} redox and balance simple cases. Learn it if you want the
complete picture (it is genuinely useful for Unit 9 electrochemistry), but
do not spend exam-prep time on it.
\end{apcnote}

\begin{exitticket}
In \ce{2 H2 + O2 -> 2 H2O}, identify the oxidizing agent and justify with
oxidation states. \answerlines[2]{\ce{O2} --- oxygen goes from 0 to $-2$
(gains electrons, is reduced), so it is the oxidizing agent. Hydrogen goes
$0 \to +1$ and is the reducing agent.}
\end{exitticket}

\end{document}
